\documentclass{article}
\usepackage[utf8]{inputenc}
\usepackage[russian]{babel}
\usepackage{amsmath}
\usepackage{amssymb}
\usepackage{mathtools}
\usepackage{amsthm}
\usepackage[a4paper, left=2cm, right=2cm, top=2cm, bottom=2cm]{geometry}

\theoremstyle{plain}
\newtheorem{theorem}{Теорема}
\newtheorem{lemma}[theorem]{Лемма}
\newtheorem{corollary}[theorem]{Следствие}

\theoremstyle{definition}
\newtheorem{definition}{Определение}
\newtheorem{example}{Пример}

\theoremstyle{remark}
\newtheorem{remark}{Замечание}

\setlength{\parindent}{1.5em}

\title{Билет 1. Неявные функции одной переменной. Теоремы о существовании и
дифференцируемости.}
\author{}
\date{}

\begin{document}

\maketitle

\section*{Определение и условие существования неявной функции}
Рассматриваем функцию двух переменных $F(x, y)$. Будем считать, что она задана
на некотором множестве $D \subset \mathbb{E}^2$.

Рассматриваем уравнение $F(x, y) = 0$. Рассматриваем задачу получения функции
$y = f(x)$.

Выбираем $\bar{x}$, подставляем в уравнение, получаем $F(\bar{x}, y) = 0$. Решаем его относительно
$y$. Находим решение $\bar{y}$. Тогда данное правило определяет $y$ как функцию от $x$.

\textbf{Более строгая формулировка}
Предпологаем, что у нас есть некоторое множество $G$ в $\mathbb{E}^1$. На этом множестве
определена функция $f(x)$. Выполнены 2 условия:
\begin{enumerate}
    \item $\forall x \in G \quad (x, f(x)) \in D$;
    \item $\forall x \in G \quad F(x, f(x)) = 0$.
\end{enumerate}
Тогда говорят, что $y=f(x)$ -- неявная функция, определяемая уравнением $F(x, y) = 0$.

\textbf{Задача поиска обратной функции как частный случай задачи о неявной функции}
Мы рассматривали функцию $x = \varphi(y)$, и хотели найти функцию $y = g(x)$ такую, что
$g(\varphi(y)) = y$. Такая задача равносильна поиску неявной функции, определяемой уравнением
$F(x, y) = \varphi(y) - x = 0$. То есть задача об обратной функции является частным случаем задачи о неявной функции.

\textbf{Условие существования неявной функции}
Рассматриваем некоторую точку $(x_0, y_0)$ и рассматриваем прямоугольник:
\[\Pi = \{(x, y) : |x - x_0| < a, |y - y_0| < b\} \subset D.\] (2)
$a, b$ -- некоторые вещественные положительные числа.

Считаем, что этот прямоугольник лежит в области определения функции $F(x, y)$.
Будем говорить, что уравнение $F(x, y) = 0$ в прямоугольнике (2) однозначным
образом задает неявную функцию $y = f(x)$, если:
\begin{enumerate}
    \item $\forall x \in (x_0 - a, x_0 + a) \quad \exists! y \in (y_0 - b, y_0 + b) : F(x, y) = 0$;
    \item множество $G = (x_0 - a, x_0 + a)$ является областью определения функции $f(x)$.
\end{enumerate}

\section*{Теорема о существовании неявной функции}
\begin{theorem}
    Пусть функция $F(x, y)$ определена в прямоугольнике:
    \[\{(x, y) : |x - x_0| < c, |y - y_0| < d\} \subset D.\] (3)
$c, d$ -- некоторые вещественные положительные числа.

    Пусть выполнены следующие условия:
    \begin{enumerate}
        \item $F(x_0, y_0) = 0$;
        \item $F(x, y)$ непрерывна в прямоугольнике (3);
        \item $\forall \bar{x} \in (x_0 - c, x_0 + c)$ выполнено условие: 
        $F(\bar{x}, y)$ строго монотонна по $y$ на $(y_0 - d, y_0 + d)$.;
    \end{enumerate}
    Тогда $\exists$ прямоугольник вида (2), в котором уравнение $F(x, y) = 0$ однозначным образом
    задает неявную функцию $y = f(x)$. Эта функция непрерывна на промежутке
    $(x_0 - a, x_0 + a)$. И удовлетворяет условию $f(x_0) = y_0$.
\end{theorem}

\begin{proof}[Доказательство]
Сначала докажем существование функции $y = f(x)$ и разбираемся как найти
прямоугольник (2).
Потом напишу.
Доказательство непрерывности этой функции.
Тоже потом. (нет, не напишу)
\end{proof}

\section*{Дифференцируемость неявной функции}
Рассматриваем уравнение $F(x, y) = 0$ (1), рассматриваем соответствующий
прямоугольник (2), $c, d$ -- некоторые вещественные положительные числа.
\begin{theorem}
    Пусть функция $F(x, y)$ определена в прямоугольнике (2) и выполнены следующие условия:
    \begin{enumerate}
        \item $F(x_0, y_0) = 0$;
        \item $F(x, y)$ непрерывна в прямоугольнике (2);
        \item $F_y'(x, y), F_x'(x, y)$ существуют и непрерывны в прямоугольнике (2);
        \item $F_y'(x_0, y_0) \neq 0$.
    \end{enumerate}
    Тогда существует прямоугольник:
    \[\{(x, y) : |x - x_0| < a, |y - y_0| < b\} \subset D. (3)\]
    в котором уравнение (1) однозначным образом задает неявную функцию
    $y = f(x)$. Эта функция непрерывна и непрерывно дифференцируема на промежутке
    $(x_0 - a, x_0 + a)$.
\end{theorem}

\begin{proof}[Доказательство]
    У нас есть условие 4. Кроме того, эта производная непрерывна, значит
    в некоторой окрестности точки $(x_0, y_0)$ она не равна нулю.
    Значит, найдется прямоугольник вида:
    \[\{(x, y) : |x - x_0| < \bar{c}, |y - y_0| < \bar{d}\} \subset (2). (\bar{2})\]
    такой, что $\forall (x, y)$ из этого прямоугольника выполнено $F_y'(x, y) \neq 0$.
    Тогда $F(\bar{x}, y)$ будет строго монотонна как функция от $y$. Таким образом,
    выполнены условия теоремы о существовании неявной функции. Следовательно, применяя
    эту теорему, получаем существование прямоугольника (3), в котором уравнение (1)
    однозначным образом задает неявную функцию $y = f(x)$.

    \textbf{Докажем дифференцируемость этой функции.}

    Берем $\bar{x} \in (x_0 - a, x_0 + a)$. Рассматриваем приращение $\Delta x$ так,
    что $\bar{x} + \Delta x \in (x_0 - a, x_0 + a)$. Данному приращению соответствует
    приращение функции $\Delta y = f(\bar{x} + \Delta x) - f(\bar{x})$.

    Рассматриваем предел:
    \[\lim_{\Delta x \to 0} \frac{\Delta y}{\Delta x} = f'(\bar{x})\]

    Необходимо доказать существование этого предела. Чтобы это сделать, отметим,
    что $f(x)$ -- неявная функция, определяемая уравнением (1). Значит, на открытом
    промежутке $(x_0 - a, x_0 + a)$ выполнено равенство:
    \[F(x, f(x)) = 0.\]
    Берем разность:
    \[F(\bar{x} + \Delta x, f(\bar{x} + \Delta x)) - F(\bar{x}, f(\bar{x})) = 0.\]
    С одном стороны, это равно нулю. С другой стороны, эта разность представима в виде:
    \[F_x'(\bar{x}, \bar{y})\Delta x + F_y'(\bar{x}, \bar{y}) \Delta y + \alpha \Delta x + \beta \Delta y\]
    Тогда, получим:
    \[\frac{\Delta y}{\Delta x} = -\frac{F_x'(\bar{x}, \bar{y}) + \alpha}{F_y'(\bar{x}, \bar{y}) + \beta}.\]
    Переходим к пределу при $\Delta x \to 0$. При этом $\Delta y \to 0$.
    Тогда $\alpha, \beta \to 0$. Следовательно, получаем:
    \[f'(\bar{x}) = -\frac{F_x'(\bar{x}, f(\bar{x}))}{F_y'(\bar{x}, f(\bar{x}))}.\]
    Теорема доказана.
\end{proof}

\end{document}